\section{把 LLM scaling recipe 迁移到多模态}

前两章建立了接口，本章讨论训练路线。课堂不是从零发明一套 multimodal science，而是把 LLM 的成功经验拆成四项可检验假设：instruction following 能否跨 modality、reasoning 能否利用视觉/音频上下文、规模扩大是否持续改善、sparse architecture 能否控制成本。每一项都可能成立，也都可能因 modality 的高维、冗余和目标差异而失效。

\subsection{四个迁移支柱}

本节把上一章的统一序列接口推进到训练路线：先明确哪些 LLM 经验值得迁移，再检查每一项在高维 modality 上可能出现的失效条件。读图时应按行为能力、规模效应和基础设施三层理解，而不是把四列当成已完成清单。

\lecturefigure{slide-019.jpg}{从 LLM paradigm 迁移出的四项 multimodal scaling 原则}{官方 deck 第 19 页；课堂 00:09:16--00:11:33}

\begin{knowledgebox}{读图：按从“行为”到“基础设施”的顺序阅读}
第一列是 prompting 与 instruction following，解决用户如何控制模型；第二列是 multimodal planning，解决模型如何在不同信息源间组织中间步骤；第三列是 data/model scaling，提出性能是否随规模连续改善；第四列是 architectural scaling，用 MoE、attention sparsity 或 modality specialization 降低 active compute。前两列主要改变 behavior，后两列决定训练和推理能否负担。
\end{knowledgebox}

\begin{importantbox}{Scaling hypothesis，而不是 scaling guarantee}
若增加数据或参数得到性能曲线 $Q(N,D)$，语言模型经验提示 $\partial Q/\partial N>0$、$\partial Q/\partial D>0$ 在一段区间内成立。但多模态数据的质量、同步误差、重复帧、缺失标注和 loss mismatch 会改变曲线。真正要验证的是 matched-compute scaling，而不是只展示更大模型的最终分数。
\end{importantbox}

\teachervoice{00:09:16--00:11:33，老师把 instruction following、planning、scale 与 architecture sparsity 视为从 LLM 迁移的原则。这里最重要的语气是“approach”：这些是研究路线，不是已经解决的工程清单。}

\subsection{两条问题轴组织全讲}

前一节列出 scaling recipe，本节把它收束成两个可检验问题：非文本输出怎样生成，以及联合训练是否产生 capability transfer。后续模型的差异都可以回到这两条轴上比较；读者还应分别记录 representation、loss 与 parameter sharing，避免用“统一模型”掩盖不同实验条件。

\lecturefigure{slide-021.jpg}{Omni model 的两条核心研究问题}{官方 deck 第 21 页；课堂 00:11:37--00:12:18}

\begin{knowledgebox}{读图：generation 与 transfer 是两个独立问题}
问题一问“怎样生成非文本 modality”，关注 representation 与 loss；问题二问“understanding 与 generation 是否正迁移”，关注共享 backbone 是否产生协同。一个模型可以解决第一问却在第二问上出现 negative transfer。Chameleon 与 Transfusion主要探索第一问；MoT 同时处理 capacity conflict；讲座后半段才正面讨论第二问。
\end{knowledgebox}

接下来可把模型设计写成五元组：
$$
\mathcal{M}=(R, O, L, P, C),
$$
其中 $R$ 是 representation，$O$ 是 ordering/sequence interface，$L$ 是 losses，$P$ 是 parameter sharing/routing，$C$ 是 compute and capacity allocation。不同论文并非互相替代，而是在五个维度上做不同选择。

\begin{knowledgebox}{阅读后续论文的统一问题表}
看到一个新 multimodal model 时，先问：输入和输出各是什么 modality？图像是离散 code 还是 continuous latent？跨 modality 是否 autoregressive？每种输出用什么 loss？attention/FFN 哪些共享？路由是 deterministic 还是 learned？比较时 active parameters、tokens、FLOPs 和数据是否匹配？如果这些问题没有答案，单看“统一模型”四个字没有足够技术含义。
\end{knowledgebox}

\subsection{本章小结}

多模态 scaling 的目标不是把所有信号强行塞进同一个 vocabulary，而是保留 sequence model 的统一控制面，同时允许 modality-specific representation、objective 与 sparse capacity。后面的三个模型族分别回答：完全离散化能走多远？连续生成目标如何与语言模型共存？参数共享何时应该被确定性稀疏化？

